\chapter{Data}
In our empirical analysis, we combined different data sources to derive a panel of subnational geographical units for the period from 2009 to 2019. Data on conflict event are drawn from the "Armed Conflict and Location Event" (ACLED) which cover conflict event from 1999 to 2024. The dataset provides the exact location of events, the type of event, and other characteristics related to the conflict. \cite{raleigh2010introducing}
In constructing the dataset, we focused on consumption-related conflict events rather than rapacity-based events. Both theoretical and empirical considerations guide this choice. Trade openness influences a country's income, usually in a positive way. However, globalisation has not always reduced inequalities in emerging countries, causing unequal distribution. 
As stated by \textcite{dube2013commodity}, an increase in income can reduce conflict, but on the other hand, higher income means there is more to fight over these resources. The existence of two opposing influences, opportunity cost and rapacity, suggests that some shocks may mitigate conflict, whereas other shocks exacerbate it, depending on the relative strength of the two effects.
The study of \textcite{amodio2024trade} shows that exposure to exports increasess simultaneously economic activity and political violence. Meanwhile, \textcite{mcguirk2020economic} have shown that, in the event of positive price shocks, producers' revenues increase but consumers' real incomes decrease, thus enhancing the probability of protests or demonstrations. Consumption-related mechanisms are more directly influenced by exogenous trade shocks, while conflicts over resources follow different dynamics, often linked to the strategies of armed groups. From an empirical point of view, limiting oneself to consumption-related events also mitigates concerns about endogeneity: resource predation can in itself influence local trade flows, risking to reflect both cause and effect simultaneously.

The final variable therefore measures the annual number of consumption-related conflict events for ADM2, filtered to retain only: (i) subcategories of events not related to predation, such as protests (peaceful or violent), demonstrations, arrests, agreements or territorial changes not directly linked to resource predation. (ii) Events reported in narrative notes as disputes over consumer goods or livelihoods (keywords such as land, livestock, pasture, farm, grazing, oil, gold, looting).
Events clearly associated with predatory motives (e.g., looting, armed clashes over resource extraction, kidnappings for ransom) are excluded.

After filtering, we limited the dataset to events occurring in coastal regions, then projected them into a consistent spatial reference system and overlaid them on Nigeria's second-level administrative units (ADM2). For each ADM2-year, the count of non-predatory conflict events is available, which serves as the outcome variable in my regression analysis.